% !TEX root = ../notes_missingsemester.tex
% Chapter 0: Projektbericht — Course Work Requirements
%%%%%%%%%%%%%%%%%%%%%%%%%%%%%%%%%%%%%%%%%%%%%%%%%%%%%%%%%%%%%%%%%%%%%%

\index{Projektbericht}\index{course work}
\chapter{Projektbericht}
\printchapterversion{00_projektbericht}
\newthought{Real artists ship. - Steve Jobs}

\section{The Why}\index{assessment}\index{project report}

A lecture that teaches tools and workflows must also assess tools and workflows.
Written exams test recall; a project report tests whether you can actually build 
and deploy something.
Enter the Projektbericht.

\marginnote{\newthought{The Projektbericht as in the Pr\"ufungsordnung.} The official frame for the Projektbericht is set by the DHBW Studien- und Pr\"ufungsordnung 
Wirtschaft, available 
\href{https://www.ravensburg.dhbw.de/fileadmin/user_upload/Dokumente/Amtliche_Bekanntmachungen/2025/27_2025_Bekanntmachung_DHBW_StuPrO_Wirtschaft_inkl._2._AES.pdf}{here}:

\textquotedblleft Die zu pr\"ufende Person soll zeigen, dass sie in der Lage ist, Projekte oder Studien mit wissenschaftlicher 
oder praktischer Problemstellung selbstst\"andig zu bearbeiten sowie deren Ergebnisse schriftlich zu dokumentieren. Der 
Projektbericht soll je zu pr\"ufender Person in Modulen mit f\"unf bis sechs ECTS-LP zehn bis zw\"olf Seiten, (bei uns entsprechend 5 Seiten) \[\dots\] umfassen.\textquotedblright}


\newthought{Pick any block from the course} and build a meaningful extension onto it.
The extension should add functionality, improve reliability, or strengthen security
in a way that goes beyond what the lecture covered.
The choice is yours, but each extension must be unique across the class:
submit your proposal for approval on a first come, first served basis with me.
If your proposal strikes home with me, we can also discuss small full stack 
projects, which live outside the scope of the lecture material.


\marginnote{\newthought{How to write the report.} A practical guide for writing reports and theses lives \href{https://pages.schutera.com/}{here}.

\medskip
\null\hfill\qrcode[height=1.0cm]{https://pages.schutera.com/}
\medskip

For deeper druidic knowledge on writing a thesis, consult \href{https://pages.schutera.com/HowToThesis/notes.pdf}{here}.

\medskip
\null\hfill\qrcode[height=1.0cm]{https://pages.schutera.com/HowToThesis/notes.pdf}
\medskip}

\newthought{Submit a written report as PDF} and a link to your GitHub repository. The report must cover:
\begin{itemize}
    \item Which block you chose and what it covers in the course.
    \item What you built on top of it, which methods, tools and techstack you used, and why it adds value.
    \item Open questions and problems you encountered and how you resolved them.
    \item What works, what remains incomplete, and what would you do next.
\end{itemize}

\newpage
\newthought{Evaluation and its Five Dimensions}.
Depth of extension asks how far the extension reaches beyond the lecture block,
in functionality, originality, method or combination.
Technical execution asks for a working repository, clean and reproducible, software engineering best practices and 
supportive of your method.
Understanding and structure shows in methods and your writing, design choices justified,
and trade-offs grasped and motivated, carried by a clear outline and a line of argument 
the reader can follow.
Reflection and rigor names problems and their resolutions, states gaps and limitations of the contribution honestly,
and makes next steps concrete.
Writing and format is concise, structured, and well argued prose
inside a correct formal frame: roughly five pages, an overview figure, a clear IMRD structure, and the reader
is always in mind.

\marginnote{\newthought{Looking for inspiration?}
Explore what tools modern teams actually use \href{https://stackshare.io}{here}
or browse the ThoughtWorks Technology Radar \href{https://www.thoughtworks.com/radar/tools}{here}.
Pick something that excites you; the best reports come from genuine curiosity.}

\marginnote{\newthought{Musterl\"osungen} live \href{https://material.schutera.com/lecturenotes/notes_missingsemester/musterloesungen/2026_MusterEins.pdf}{here}.
Use them to calibrate scope and structure, then write your own.

\medskip
\null\hfill\qrcode[height=1.0cm]{https://material.schutera.com/lecturenotes/notes_missingsemester/musterloesungen/2026_MusterEins.pdf}
\medskip}
